\section{الفصل~\ref{sec:sensors}. مداخل الآلة}

بنية المستشعرات مقيسة مباشرة، بتسجيل تيار الخلايا المفردة واهتزازات القوقعة وبعدّ خلايا الشبكية؛ وحد الحساسية وصيغة ضجيج الطلقات يتبعان من الفيزياء؛ أما أن شفرات المستشعرات مضبوطة على أكبر قدر من المعلومات ففرضية لها سند قوي.

{\footnotesize
\setlength{\tabcolsep}{3pt}\renewcommand{\arraystretch}{1.15}
\begin{longtable}{>{\raggedright}p{0.5cm}>{\raggedright}p{4.0cm}>{\raggedright}p{5.6cm}>{\raggedright}p{2.3cm}>{\raggedright\arraybackslash}p{1.7cm}}
\toprule
م & القضية & التجربة وما قيس & المصدر & الحالة \\
\midrule\endfirsthead
\toprule
م & القضية & التجربة وما قيس & المصدر & الحالة \\
\midrule\endhead
1 & المستشعر محوِّل: المستقبِل يعطي تيارًا، ومخرج الخلايا الحسية في العين والأذن غلوتامات على ناقل \nt{GLUT}؛ والخلايا الأولى لا تُصدر نبضات (الشكل~\ref{fig:transducer}) & العصي والمخاريط لدى السلحفاة تطلق حمضًا أمينيًا مثيرًا: تلتقطه قنوات غلوتاماتية على رقعة غشاء مفصولة. الخلايا الشعرية الداخلية في قوقعة الجرذ: التيارات في نهايات الألياف العصبية تمر عبر مستقبِلات \foreignlanguage{english}{AMPA} الغلوتاماتية، ويزداد تردد التحرير مع إزالة الاستقطاب المتدرجة للخلية حتى $150\,\text{Hz}$ & \cite{CopenhagenJahr1989,GlowatzkiFuchs2002} & مُقاس \\
2 & مستشعر الضوء في الظلام يستجيب لفوتون واحد بتيار نحو بيكوأمبير & قطب ماصّ على القطعة الخارجية لعصية العلجوم: الاستجابات للومضات الخافتة مكمّاة، والحدث $\approx1\,\text{pA}$ ($5\%$ من التيار المظلم) بانتشار $0.2\,\text{pA}$ وكفاءة $0.5$. والإنسان يُبلغ عن وصول فوتون واحد أكثر مما تقتضيه المصادفة & \cite{Baylor1979,Tinsley2016} & مُقاس \\
3 & مستشعر الصوت يلاحظ إزاحة أقل من نانومتر & طريقة موسباور، الغشاء القاعدي لقوقعة خنزير غينيا عند موضع $16\text{--}19\,\text{kHz}$: عند عتبة استجابة العصب السمعي سرعة الغشاء نحو $0.04\,\text{mm/s}$، أي إزاحة $\approx0.35\,\text{nm}$ (التحويل من عندنا). المقيس هو الغشاء لا حزمة المستشعر & \cite{Sellick1982,Hudspeth1989} & مُقاس \\
4 & الدقة في الظلام يحدّها ضجيج الطلقات للفوتونات~\eqref{eq:photonnoise}؛ وفي الضوء الخافت يطول التجميع & الصيغة تتبع من إحصاء بواسون. في العصية يطابق الضجيجُ في الضوء الخافت مجموعَ أحداث كمومية مستقلة؛ وعلى خلفية ساطعة تكون الاستجابة للومضة أصغر وأقصر. أما رقم ”أعشار الثانية“ فلم يُختبر هنا على حدة & \cite{Baylor1979} & نظرية \\
5 & مدى الدخل عشر رتب للضوء واثنتا عشرة للصوت، أما تردد النبضات فأقل من ثلاث رتب & للصوت: تزداد استجابة الغشاء القاعدي عند التردد المميِّز بميل يصل إلى $0.2\,\text{dB/dB}$ في النطاق $40\text{--}80\,\text{dB}$، والانضغاط ظاهر فوق $100\,\text{dB}$ أيضًا. أما أعداد الرتب نفسها فتقديرات معيارية، بلا مصدر في الفصل & \cite{Ruggero1997} & مُقاس \\
6 & استجابة المستشعر~\eqref{eq:nakarushton}، و$\sigma$ تتبع متوسط السطوع فتزيح المنحنى على المحور اللوغاريتمي (الشكل~\ref{fig:adaptation}) & طُبّقت الصيغة أول مرة على كمونات \foreignlanguage{english}{S} في شبكية السمك. مخاريط السلحفاة: الاستجابات تخضع لـ$V=I V_m/(I+\sigma)$ وتغطي $\sim3.5$ رتبة من السطوع؛ وبعد $2\text{--}3$ دقائق من الخلفية ينزاح المنحنى أفقيًا محافظًا على عرضه. الصلة بالقسمة على النشاط الكلي مراجعة & \cite{NakaRushton1966,NormannPerlman1979,Carandini2012} & مُقاس \\
7 & المستشعرات تنقل التغيرات، وشفراتها مضبوطة على أكبر قدر من المعلومات لكل نبضة (القضية~\ref{prop:sensors}) & المبدأ مقترح نظريًا. وأقوى دليل: لدى الذبابة تطابق خاصية ”التباين $\leftarrow$ الاستجابة“ لعصبونات الطبقة الأولى التوزيعَ التراكمي لتباينات المشاهد الطبيعية، وبذلك تُبلغ أكبر كمية من المعلومات & \cite{Barlow1961,Laughlin1981} & فرضية \\
8 & مستشعرات التشغيل والإطفاء شفرة سلكين؛ وكاميرا الأحداث تكررها في السيليكون & مراجعة للتجارب: قناتا التشغيل والإطفاء في الشبكية تنفصلان منذ المشبك الأول ويمكن إيقاف كل منهما على حدة. الكاميرا: $128\times128$ بكسل بلا إطارات، مدى $120\,\text{dB}$، وتأخير $15\,\text{\textmu s}$ & \cite{Schiller1992,Lichtsteiner2008,Gallego2022} & مُقاس \\
9 & في العين $\sim10^8$ مستشعر ضوء و$\sim10^6$ عصبون مخرج: ضغط بنحو $\sim100$ مرة & عدّ على شبكيات بشرية كاملة: في المتوسط $92$ مليون عصية و$4.6$ مليون مخروط؛ وخلايا المخرج (العقدية) من $0.7$ إلى $1.5$ مليون في عيون مختلفة & \cite{Curcio1990,CurcioAllen1990} & مُقاس \\
10 & لدى الفأر أكثر من ثلاثين قناة مخرج للعين & الفأر: تصوير كالسيومي ثنائي الفوتون لأكثر من $11\,000$ خلية مخرج وتجميع الاستجابات، فظهر أكثر من $30$ نوعًا وظيفيًا بفارق واضح. لدى الإنسان لم يُقَس العدد & \cite{Baden2016} & مُقاس \\
11 & عين خنزير غينيا تنقل $\sim875\,000$ بت في الثانية؛ والتحويل إلى عين الإنسان يعطي $\sim10^7$ بت في الثانية & خنزير غينيا، مصفوفة متعددة الأقطاب، حركة في مشاهد طبيعية: $6\text{--}13$ بت في الثانية لكل خلية، و$\sim875\,000$ بت في الثانية لنحو $\sim10^5$ خلية مخرج. وللإنسان ($\sim10^6$ خلية) $\sim10^7$ بت في الثانية، وهو تحويل المؤلفين & \cite{Koch2006} & مُقاس \\
12 & الأذن بنك مرشحات تمرير النطاق بخريطة ترددات شبه لوغاريتمية (الشكل~\ref{fig:filterbank}) & موضع أكبر اهتزاز للغشاء القاعدي يتوقف على التردد؛ ودالة ”التردد--الموضع“ شبه أُسية، وتصف بشكل واحد قواقع الإنسان والحيوانات الحية & \cite{Greenwood1990,Robles2001} & مُقاس \\
13 & المرنانات فعّالة: بعض المستشعرات يضخّم الاهتزازات الضعيفة آلاف المرات في السعة ($66\text{--}76\,\text{dB}$)، ومع ازدياد الشدة ينخفض الكسب & قياس الاهتزاز بالليزر عند قاعدة قوقعة الشنشيلة: للنغمات الضعيفة عند التردد المميِّز كسب $66\text{--}76\,\text{dB}$ نسبةً إلى الركاب؛ والموت يخفض الحساسية بـ$60\text{--}81\,\text{dB}$ (بـ$10^3\text{--}10^4$ مرة في السعة) ويزيل الانضغاط. مصدر القوة الخلايا الشعرية الخارجية & \cite{Ruggero1997,Robles2001} & مُقاس \\
14 & مضخّم الأذن موضوع عند حافة الإثارة الذاتية & حساب: قرب نقطة تفرّع هوبف لا مفر من انضغاط المدى والتوليف الحاد والنغمات التركيبية، وهي كل اللاخطيات المعروفة في السمع. أما أن القوقعة مضبوطة هكذا بالذات فلم يُبرهن مباشرة & \cite{Eguiluz2000} & فرضية \\
15 & في الترددات المنخفضة تسير النبضات في طور مع الصوت (لدى خنزير غينيا يضعف التقيد بدءًا من $\sim600\,\text{Hz}$ ويبقى ملحوظًا حتى $\sim3.5\,\text{kHz}$)، وبها يتحدد الاتجاه؛ وفي العالية تبقى شفرة الموضع & العصب السمعي لخنزير غينيا: يبدأ التقيد بالطور بالانخفاض قرب $600\,\text{Hz}$ ويختفي فوق $3.5\,\text{kHz}$؛ ولدى القط والقرد الحد أعلى بنحو أوكتاف. فرق زمن الوصول إلى الأذنين مراجعة & \cite{PalmerRussell1986,Grothe2010} & مُقاس \\
16 & بقية المستشعرات (الجدول~\ref{tab:sensors}): للشم $\sim400$ نوع من المستقبِلات، نوع لكل عصبون، وشفرة تركيبية؛ والذوق جزئيًا عبر \foreignlanguage{english}{ATP} والسيروتونين؛ واللمس مستشعرات سريعة التعوّد وبطيئته؛ والدفء والبرد منفصلان & الفأر: المستقبِل الواحد يتعرف على روائح كثيرة، والرائحة الواحدة على مستقبِلات كثيرة، والروائح المختلفة على مجموعات مختلفة. بلا مستقبِلات \foreignlanguage{english}{ATP} لا يستجيب عصب الذوق؛ والبراعم الذوقية تطلق السيروتونين. تسجيل ألياف مفردة من جلد يد الإنسان: أربعة أنواع بحسب سرعة التعوّد. وقناة البرد مختلفة عن قنوات الحرارة & \cite{Mombaerts2004,Malnic1999,Finger2005,Huang2005,JohanssonVallbo1979,McKemy2002} & مُقاس \\
\bottomrule
\end{longtable}}
